\section{Consequences of localization and parity}
\label{sec:consequences}

The direct parity proof is complete. Localization next bounds
complementary-output probabilities, and a construction of Gretta, Gupta,
and Joshi~\cite{GGJ2026} gives a Dicke-state obstruction. We also derive
scalar-output Fourier concentration directly from localization.
This concentration rules out high average accuracy for majority,
while Xu and Li's reductions~\cite{XuLi2026Fanout} give worst-case
obstructions for symmetric functions. Throughout, discarded registers
remain unrestricted.

\subsection{State preparation}
\label{sec:state-preparation}

Gretta, Gupta, and Joshi~\cite{GGJ2026} connect complementary-output
probabilities and Dicke-state preparation to parity. Here localization
bounds the complementary probabilities directly; their block construction
then gives the Dicke-state consequence.
For an $n$-qubit density operator $\rho$, write
\[
 p_y=\bra{y}\rho\ket{y},\qquad
 \mathcal F_n(\rho)=2\sum_{y\in\{0,1\}^n}p_y p_{\bar y},
\]
where $\bar y$ is the bitwise complement of $y$.
This is the \emph{felinity} of $\rho$~\cite[Definition~2.1]{GGJ2026}.
Also write
\[
 \ket{D_k^n}=\binom{n}{k}^{-1/2}\sum_{|y|=k}\ket{y},
 \qquad
 \operatorname{TD}(\rho,\sigma)=\tfrac12\norm{\rho-\sigma}_1
\]
for the weight-$k$ Dicke state and trace distance, respectively.

\begin{corollary}[State-preparation obstructions]
\label{cor:state-preparation}
Fix a nonnegative integer $d$, $c\ge1$, $a>0$, and $0<\delta<1$.
For all sufficiently large $n$, every $n$-qubit reduced state $\rho$
prepared from zero inputs by an allowed depth-at-most-$d$ circuit on
at most $n^c$ total qubits satisfies
\[
 \mathcal F_n(\rho)<n^{-a}.
\]
Moreover, for every integer $k$ with $n^\delta\le k\le n/2$,
\[
 \operatorname{TD}(\rho,\proj{D_k^n})>\frac1{80k}.
\]
The remaining qubits may be discarded in arbitrary final states.
\end{corollary}

\begin{proof}
Retain all $N$ preparation qubits, including those eventually discarded,
and write $\ket\psi=U\ket{0^N}$ and $Q=\proj\psi$. Thus $Q$ is the
conjugate by $U\in\cU_d(N)$ of a product projection. For an output string
$y$, let $P_y$ fix the $n$ output qubits to $y$ and act identically on the
remaining register, and let $M_y$ count mismatches from $y$ on the outputs.
Then $P_y=[M_y=0]$ and $\norm{P_y\psi}^2=p_y$. Fix $0<s<t<1/c$.
For sufficiently large $n$, $N^t<n$, so the complementary-string
projection satisfies $P_{\bar y}\le[M_y\ge N^t]$. Localization gives
\[
 \sqrt{p_y p_{\bar y}}
 =\norm{P_{\bar y}QP_y}\le e^{-N^s}
 \qquad\text{for every }y.
\]
The threshold is uniform in $y$. Cauchy--Schwarz gives
$\sum_y\sqrt{p_y p_{\bar y}}\le1$, and hence
\[
 \mathcal F_n(\rho)
 \le2e^{-N^s}\sum_y\sqrt{p_y p_{\bar y}}
 \le2e^{-N^s}<n^{-a}
\]
eventually. This argument uses the pure state on the complete register,
so it applies to every reduced output state $\rho$.

For the second assertion, we give the block construction underlying
\cite[Lemma~4.13 and Theorem~4.14]{GGJ2026}, including the estimates
needed at the stated trace distance. Put
\[
 \ell=\lfloor k/\log k\rfloor,\qquad p=k/n,\qquad
 \ket v=\sqrt{1-p}\ket0+\sqrt p\ket1.
\]
Partition the $n$ sites into $\ell$ blocks $B_j$ whose sizes $b_j$ differ
by at most one. Append one zero qubit $t_j$ per block, and on $B_jt_j$ apply
\[
 (I-A_j)\otimes I+A_j\otimes X,
 \qquad A_j=(\proj v)^{\otimes b_j}.
\]
Each gate equals $I-2A_j\otimes\proj{-}$, so all these gates form one
disjoint product-reflection layer, implemented by parallel one-qubit basis
changes, one Toffoli layer, and the inverse basis changes.
Let $\sigma_*$ be the resulting state of the $\ell$ new qubits when the
input is $\ket{D_k^n}$. Its all-one probability is
\[
 r_1=|\langle v^{\otimes n}\mid D_k^n\rangle|^2
     =\binom nk p^k(1-p)^{n-k}.
\]
For a binomial random variable $B$ with parameters $n,p$, the integer $k$
is a mode and $\operatorname{Var}(B)\le k$. Chebyshev's inequality gives
\[
 (4\sqrt k+1)r_1\ge\Pr(|B-k|<2\sqrt k)\ge\frac34,
 \qquad r_1\ge\left(\frac3{16}-o(1)\right)k^{-1/2},
\]
since the displayed interval contains at most $4\sqrt k+1$ integers.
This estimate is uniform in $n\ge2k$. The all-zero probability is
$r_0=\norm{\Pi\ket{D_k^n}}^2$, where $\Pi=\prod_j(I-A_j)$.
To bound it, use the unit vector
\[
 \ket\nu=(\sqrt{1-p}\ket0-\sqrt p\ket1)^{\otimes n}.
\]
Since $b_j\ge n/\ell-1$, $p\le1/2$, and $k/\ell\ge\log k$,
\[
 \langle\nu|A_j|\nu\rangle=(1-2p)^{2b_j}
 \le e^{-4pb_j}\le e^2k^{-4}.
\]
The commuting projections satisfy $I-\Pi\le\sum_jA_j$, so
$\norm{(I-\Pi)\nu}\le e k^{-3/2}$. Moreover,
$|\langle\nu\mid D_k^n\rangle|=\sqrt{r_1}$, because every term of the
Dicke state has weight $k$. Therefore
\[
 \sqrt{r_0}\ge\sqrt{r_1}-e k^{-3/2},\qquad
 \mathcal F_\ell(\sigma_*)\ge4r_0r_1
       \ge\frac{9/64-o(1)}k>\frac1{8k}.
\]
All errors here are uniform for $k\le n/2$ as $k\to\infty$.
The factor four counts both complementary strings in the definition
of felinity.

Apply the same block circuit to $\rho$, and call its new-qubit state
$\sigma$. Trace distance contracts under a unitary and partial trace.
Also, for any two states $\tau,\tau'$ with computational-basis
probabilities $q,q'$, the definition gives
\[
 |\mathcal F_\ell(\tau)-\mathcal F_\ell(\tau')|
 \le4\sum_y|q_y-q'_y|
 \le8\operatorname{TD}(\tau,\tau').
\]
Thus $\operatorname{TD}(\rho,\proj{D_k^n})\le1/(80k)$ would imply
\[
 \mathcal F_\ell(\sigma)>\frac1{8k}-\frac8{80k}
 =\frac1{40k}.
\]
For large $k$, $\ell\ge\sqrt k$, while $k\ge n^\delta$ gives
$n\le\ell^{2/\delta}$. The preparation of $\sigma$ therefore has fixed
depth and polynomially many total qubits as a function of $\ell$.
But $\ell^2/k\to\infty$, so its felinity exceeds $\ell^{-2}$ eventually,
contrary to the first assertion applied with target size $\ell$ and
exponent $a=2$. Since $k\ge n^\delta\to\infty$, all thresholds can be
chosen uniformly over the allowed $k$.
\end{proof}

In particular, no such family can assign inverse-polynomial
probabilities to both a string and its complement: their two terms
contribute $4p_y p_{\bar y}$ to $\mathcal F_n(\rho)$.

The polynomial bound in \Cref{cor:state-preparation} is measured in the
number $n$ of target qubits. Gretta, Gupta, and
Joshi~\cite[Theorem~1.1]{GGJ2026StatePreparation} show that every $n$-qubit
pure state can be prepared exactly at constant depth with $2^{O(n)}$
ancillas, all returned to zero.

\subsection{Scalar-output Fourier concentration}

Gretta, Gupta, and Joshi~\cite[Conjecture~1 and Corollary~1.2]{GGJ2026}
formulate a qualitative Quantum-LMN statement and establish its
equivalence with the parity lower-bound question. Our localization
theorem gives this Fourier-concentration statement directly.
For an $n$-input circuit $C$ on $N$ total qubits with output qubit $t$,
define its expected measured sign and uniform-input Fourier coefficients by
\begin{align*}
 f_C(x)&=\bra{x,0^{N-n}}C^\dagger Z_tC\ket{x,0^{N-n}}\in[-1,1],\\
 \widehat f_C(S)&=2^{-n}\sum_{x\in\{0,1\}^n}
 f_C(x)(-1)^{\sum_{i\in S}x_i}.
\end{align*}
Write $W_{\ge k}[f_C]=\sum_{|S|\ge k}|\widehat f_C(S)|^2$.
This is the absolute squared Fourier tail, not a fraction of the
total weight $\mathbb E_x f_C(x)^2\le1$.

\begin{corollary}[Qualitative scalar-output Quantum-LMN]
\label{cor:quantum-lmn}
For every fixed nonnegative integer $d$, there is a positive function
$\gamma_d$ on the positive integers, depending only on $d$, such that
$\gamma_d(k)/k^A\to\infty$ for every fixed $A>0$, and every allowed
depth-at-most-$d$ circuit $C$ on $n$ inputs and $N$ total qubits,
every designated output $t$, and every integer $1\le k\le n$ satisfy
\[
 W_{\ge k}[f_C]\le\frac{N}{\gamma_d(k)}.
\]
In particular, fix $c\ge1$, $0<\delta\le1$, and $B>0$. For all
sufficiently large $n$, uniformly over these circuits with $N\le n^c$,
\[
 W_{\ge\lceil n^\delta\rceil}[f_C]\le n^{-B}.
\]
The sufficient threshold depends only on $d,c,\delta,B$.
\end{corollary}

\begin{proof}
Fix a circuit $C$ and let $Q=C^\dagger O_0C$, where $O_0$ projects its
output qubit onto $\ket0$. For $h\in\{0,1\}^n$, let $\xi_h$ be the
Hadamard product vector from the proof of \Cref{thm:parity}, and let
$M_h$ count mismatches from its factors on the $n$ input qubits.
For $S\subseteq[n]$, write $\mathbf 1_S$ for its indicator string.
The character calculation in \eqref{eq:parity-average}, with
$\bar h$ replaced by $h\oplus\mathbf 1_S$, gives
\[
 \widehat f_C(S)
 =2\E_h\bra{\xi_{h\oplus\mathbf 1_S},0}Q\ket{\xi_h,0}
 \qquad(S\ne\varnothing),
\]
since $f_C(x)=2\bra{x,0}Q\ket{x,0}-1$.
For fixed $h$, the vectors $\ket{\xi_{h\oplus\mathbf 1_S},0}$ are
orthonormal as $S$ varies, and their $M_h$ counts are exactly $|S|$.
Jensen's and Bessel's inequalities therefore yield
\begin{align}
 W_{\ge k}[f_C]
 &\le4\E_h\sum_{|S|\ge k}
   \left|\bra{\xi_{h\oplus\mathbf 1_S},0}Q\ket{\xi_h,0}\right|^2
   \notag\\
 &\le4\sup_h\norm{[M_h\ge k]Q[M_h=0]}^2.
 \label{eq:fourier-localization}
\end{align}
For any fixed $0<\sigma<\tau$, \Cref{thm:localization} bounds this
by $4e^{-2N^\sigma}$ whenever $N$ is sufficiently large and
$k\ge N^\tau$. The threshold depends only on $d,\sigma,\tau$.

It remains to obtain one function of $k$ valid for every circuit size.
For each positive integer $k$, define
\[
 b_d(k)=\sup\frac{W_{\ge k}[f_C]}{N},
\]
where the supremum ranges over all allowed depth-at-most-$d$ circuits,
all designated outputs, and all input and total sizes $k\le n\le N$.
Parseval's identity gives $W_{\ge k}[f_C]\le\E_x f_C(x)^2\le1$,
so $0\le b_d(k)\le1/k$.
Fix $A\ge1$. If $N\ge k^A$, then $W_{\ge k}[f_C]/N\le k^{-A}$.
For the remaining sizes $k\le N<k^A$, choose fixed
$0<\sigma<\tau<1/A$. For sufficiently large $k$, we have
$N^\tau<k$, so \eqref{eq:fourier-localization} gives
\[
 \frac{W_{\ge k}[f_C]}{N}\le\frac{4e^{-2k^\sigma}}{k}.
\]
These two bounds are uniform over the supremum. Since $A$ is arbitrary,
$b_d(k)=o(k^{-B})$ for every fixed $B>0$.
Thus the positive function
\[
 \gamma_d(k)=\frac1{b_d(k)+e^{-k}}
\]
satisfies $\gamma_d(k)/k^B\to\infty$ for every $B>0$, and the definition
of $b_d$ gives $W_{\ge k}[f_C]\le N/\gamma_d(k)$.
Finally, choose $B'>(c+B)/\delta$ in this superpolynomial growth
property to obtain the asserted cutoff consequence.
\end{proof}

\subsection{Majority and symmetric functions}

The reductions of Xu and Li turn the parity lower bound into a
worst-case lower bound for majority and broader symmetric families.
For a symmetric \(f_n:\{0,1\}^n\to\{0,1\}\), let \(f_{n,i}\) be its
value on inputs of Hamming weight \(i\), and define its transition radius by
\[
 \rho(f_n)=\max_{\substack{0\le i<n\\f_{n,i}\ne f_{n,i+1}}}
                 \min\{i+1,n-i\},
\]
with \(\rho(f_n)=0\) for a constant function. Write
\(\widetilde{\deg}_{1/3}(f_n)\) for the least total degree of a real
polynomial approximating \(f_n\) pointwise within \(1/3\).
For nonconstant symmetric functions, Paturi's characterization
\cite[Theorem~1]{Paturi1992} gives
\[
 \widetilde{\deg}_{1/3}(f_n)=\Theta\!\left(\sqrt{n\rho(f_n)}\right).
\]
Thus the approximate-degree hypothesis below also forces a polynomially
large transition radius.

\begin{corollary}[Majority and symmetric functions]
\label{cor:symmetric-functions}
Let \((f_n)\) be a symmetric Boolean family satisfying either
\(\rho(f_n)\ge n^\delta\) for some fixed \(\delta>0\), or
\(\widetilde{\deg}_{1/3}(f_n)=\Omega(n^{1/2+\eta})\) for some fixed
\(\eta>0\), for all sufficiently large \(n\). For every fixed
\(\gamma>0\), no nonuniform family in the circuit model of
\Cref{thm:parity}, with fixed depth and polynomially many total qubits,
computes \(f_n\) with probability at least
\(1/2+1/(\log n)^\gamma\) on every input for all sufficiently large \(n\).
In particular, this excludes strict majority
\(\operatorname{MAJ}_n(x)=\mathbf 1\{|x|>n/2\}\), whose transition radius
is \(\lceil n/2\rceil\).
\end{corollary}

\begin{proof}
Xu and Li's Theorem~16 and Corollaries~2--3
\cite{XuLi2026Fanout} turn a family satisfying either hypothesis and the
stated success bound into an exact, clean implementation of fanout on
\(n\) targets, with constant entangling depth and polynomially many
ancillas. Their circuit convention~\cite[Section~2.3]{XuLi2026Fanout}
requires the measured output to be an initially zero ancilla; appending one fresh zero qubit and a CNOT from our designated
output preserves its measurement distribution, even with arbitrary
garbage. Their entangling depth \(D\) uses at most \(2D+1\) of our layers.
Thus the reduction stays within fixed depth and polynomial total qubits.
Hadamard conjugation on the fanout data wires gives the exact parity
unitary; a zero target and its final measurement contradict
\Cref{thm:parity}. Finite exceptional input lengths do not affect this
circuit-family contradiction.
\end{proof}

There is consequently a strict separation between the Boolean families
computed with worst-case success at least \(2/3\) in our model and in the
matched model obtained by adjoining unit-cost reversible unbounded-arity
threshold gates
\[
 \ket{x,b}\longmapsto
 \ket{x,b\mathbin{\oplus}\mathbf 1\{|x|\ge t\}},
 \qquad t\in\{0,\ldots,n\},
\]
while retaining the same depth, disjoint-support, qubit, ancilla, and
output conventions. The enlarged model contains the original one and
computes strict majority exactly with one such gate at
\(t=\lfloor n/2\rfloor+1\), whereas
Corollary~\ref{cor:symmetric-functions} excludes even worst-case success \(2/3\)
for majority in the original model.

\subsection{High average accuracy for majority}

The preceding obstruction assumes an advantage on every input. The
high-average-accuracy consequence of Gretta, Gupta, and
Joshi~\cite[Corollary~4.6]{GGJ2026} also follows directly from our
Fourier-concentration bound: majority retains a polynomial Fourier tail
at a polynomially growing cutoff.

\begin{corollary}[High-average-accuracy majority obstruction]
\label{cor:majority-average}
For every fixed \(\delta>0\), no nonuniform family \((C_n)\) in the
circuit model of \Cref{thm:parity}, with fixed depth and polynomially
many total qubits, satisfies
\[
 2^{-n}\sum_{x\in\{0,1\}^n}
 f_{C_n}(x)(-1)^{\operatorname{MAJ}_n(x)}\ge1-n^{-\delta}
\]
for all sufficiently large \(n\). Here \(f_{C_n}(x)\in[-1,1]\) is
the expected measured sign: the probability of output zero minus the
probability of output one. Equivalently, no such family computes strict
majority with average success at least \(1-\tfrac12 n^{-\delta}\),
where the average is over uniform inputs and the output measurement.
\end{corollary}

\begin{proof}
Suppose such a family exists, restrict to odd $n$, and put
$g_n(x)=(-1)^{\operatorname{MAJ}_n(x)}$ and
$\delta_0=\min\{\delta,1/4\}$. Since $|f_{C_n}|\le1$ and $g_n^2=1$,
the asserted correlation gives
\[
 \|f_{C_n}-g_n\|_2^2
 \le 2\bigl(1-\mathbb E_x f_{C_n}(x)g_n(x)\bigr)
 \le 2n^{-\delta_0},
\]
where the norm uses the uniform probability measure.

We use the following finite-size majority tail estimate:
for $\varepsilon\ge n^{-1/3}$ and sufficiently large odd $n$,
\[
 W_{\ge k}[g_n]>\varepsilon
 \quad\text{whenever $k$ is odd and }1\le k\le\frac1{7\varepsilon^2}.
\]
This is \cite[Theorem~3.5.6, p.~50]{ODonnell2003Noise}.
Set $\varepsilon=4n^{-\delta_0}$ and let $k_n$ be the largest odd
integer at most $n^{2\delta_0}/112$. These parameters satisfy the
estimate's hypotheses for all sufficiently large odd $n$.
Parseval's identity and the triangle inequality for the Fourier
coefficients at degrees at least $k_n$ now give
\[
 \sqrt{W_{\ge k_n}[f_{C_n}]}
 \ge \sqrt{W_{\ge k_n}[g_n]}-\|f_{C_n}-g_n\|_2
 > (2-\sqrt2)n^{-\delta_0/2}.
\]
On the other hand, $k_n\ge\lceil n^{\delta_0}\rceil$ eventually.
Applying \Cref{cor:quantum-lmn} at the latter cutoff, with
$B=\delta_0+1$, yields
$W_{\ge k_n}[f_{C_n}]\le n^{-(\delta_0+1)}$, a contradiction.
\end{proof}

This conclusion does not assert a constant average-case gap or exclude
\(1-1/\operatorname{polylog}n\) correlation.
